\documentclass[journal,12pt,twocolumn]{IEEEtran}
%
\usepackage{setspace}
\usepackage{gensymb}
%\doublespacing
\singlespacing

%\usepackage{graphicx}
%\usepackage{amssymb}
%\usepackage{relsize}
\usepackage[cmex10]{amsmath}
%\usepackage{amsthm}
%\interdisplaylinepenalty=2500
%\savesymbol{iint}
%\usepackage{txfonts}
%\restoresymbol{TXF}{iint}
%\usepackage{wasysym}
\usepackage{amsthm}
\usepackage{iithtlc}
\usepackage{mathrsfs}
\usepackage{txfonts}
\usepackage{stfloats}
\usepackage{bm}
\usepackage{cite}
\usepackage{cases}
\usepackage{subfig}
%\usepackage{xtab}
\usepackage{longtable}
\usepackage{multirow}
%\usepackage{algorithm}
%\usepackage{algpseudocode}
\usepackage{enumitem}
\usepackage{mathtools}
\usepackage{tikz}
\usepackage{circuitikz}
\usepackage{verbatim}
\usepackage{tfrupee}
\usepackage[breaklinks=true]{hyperref}
%\usepackage{stmaryrd}
\usepackage{tkz-euclide} % loads  TikZ and tkz-base
\usetkzobj{all}
\usepackage{listings}
    \usepackage{color}                                            %%
    \usepackage{array}                                            %%
    \usepackage{longtable}                                        %%
    \usepackage{calc}                                             %%
    \usepackage{multirow}                                         %%
    \usepackage{hhline}                                           %%
    \usepackage{ifthen}                                           %%
  %optionally (for landscape tables embedded in another document): %%
    \usepackage{lscape}     
\usepackage{multicol}
\usepackage{chngcntr}
%\usepackage{enumerate}

%\usepackage{wasysym}
%\newcounter{MYtempeqncnt}
\DeclareMathOperator*{\Res}{Res}
%\renewcommand{\baselinestretch}{2}
\renewcommand\thesection{\arabic{section}}
\renewcommand\thesubsection{\thesection.\arabic{subsection}}
\renewcommand\thesubsubsection{\thesubsection.\arabic{subsubsection}}

\renewcommand\thesectiondis{\arabic{section}}
\renewcommand\thesubsectiondis{\thesectiondis.\arabic{subsection}}
\renewcommand\thesubsubsectiondis{\thesubsectiondis.\arabic{subsubsection}}

% correct bad hyphenation here
\hyphenation{op-tical net-works semi-conduc-tor}
\def\inputGnumericTable{}                                 %%

\lstset{
%language=C,
frame=single, 
breaklines=true,
columns=fullflexible
}
%\lstset{
%language=tex,
%frame=single, 
%breaklines=true
%}

\begin{document}
%


\newtheorem{theorem}{Theorem}[section]
\newtheorem{problem}{Problem}
\newtheorem{proposition}{Proposition}[section]
\newtheorem{lemma}{Lemma}[section]
\newtheorem{corollary}[theorem]{Corollary}
\newtheorem{example}{Example}[section]
\newtheorem{definition}[problem]{Definition}
%\newtheorem{thm}{Theorem}[section] 
%\newtheorem{defn}[thm]{Definition}
%\newtheorem{algorithm}{Algorithm}[section]
%\newtheorem{cor}{Corollary}
\newcommand{\BEQA}{\begin{eqnarray}}
\newcommand{\EEQA}{\end{eqnarray}}
\newcommand{\define}{\stackrel{\triangle}{=}}

\bibliographystyle{IEEEtran}
%\bibliographystyle{ieeetr}


\providecommand{\mbf}{\mathbf}
\providecommand{\pr}[1]{\ensuremath{\Pr\left(#1\right)}}
\providecommand{\qfunc}[1]{\ensuremath{Q\left(#1\right)}}
\providecommand{\sbrak}[1]{\ensuremath{{}\left[#1\right]}}
\providecommand{\lsbrak}[1]{\ensuremath{{}\left[#1\right.}}
\providecommand{\rsbrak}[1]{\ensuremath{{}\left.#1\right]}}
\providecommand{\brak}[1]{\ensuremath{\left(#1\right)}}
\providecommand{\lbrak}[1]{\ensuremath{\left(#1\right.}}
\providecommand{\rbrak}[1]{\ensuremath{\left.#1\right)}}
\providecommand{\cbrak}[1]{\ensuremath{\left\{#1\right\}}}
\providecommand{\lcbrak}[1]{\ensuremath{\left\{#1\right.}}
\providecommand{\rcbrak}[1]{\ensuremath{\left.#1\right\}}}
\theoremstyle{remark}
\newtheorem{rem}{Remark}
\newcommand{\sgn}{\mathop{\mathrm{sgn}}}
\providecommand{\abs}[1]{\left\vert#1\right\vert}
\providecommand{\res}[1]{\Res\displaylimits_{#1}} 
\providecommand{\norm}[1]{\left\lVert#1\right\rVert}
%\providecommand{\norm}[1]{\lVert#1\rVert}
\providecommand{\mtx}[1]{\mathbf{#1}}
\providecommand{\mean}[1]{E\left[ #1 \right]}
\providecommand{\fourier}{\overset{\mathcal{F}}{ \rightleftharpoons}}
%\providecommand{\hilbert}{\overset{\mathcal{H}}{ \rightleftharpoons}}
\providecommand{\system}{\overset{\mathcal{H}}{ \longleftrightarrow}}
	%\newcommand{\solution}[2]{\textbf{Solution:}{#1}}
\newcommand{\solution}{\noindent \textbf{Solution: }}
\newcommand{\cosec}{\,\text{cosec}\,}
\providecommand{\dec}[2]{\ensuremath{\overset{#1}{\underset{#2}{\gtrless}}}}
\newcommand{\myvec}[1]{\ensuremath{\begin{pmatrix}#1\end{pmatrix}}}
\newcommand{\mydet}[1]{\ensuremath{\begin{vmatrix}#1\end{vmatrix}}}
%\numberwithin{equation}{section}
\numberwithin{equation}{subsection}
%\numberwithin{problem}{section}
%\numberwithin{definition}{section}
\makeatletter
\@addtoreset{figure}{problem}
\makeatother

\let\StandardTheFigure\thefigure
\let\vec\mathbf
%\renewcommand{\thefigure}{\theproblem.\arabic{figure}}
\renewcommand{\thefigure}{\theproblem}
%\setlist[enumerate,1]{before=\renewcommand\theequation{\theenumi.\arabic{equation}}
%\counterwithin{equation}{enumi}


%\renewcommand{\theequation}{\arabic{subsection}.\arabic{equation}}

\def\putbox#1#2#3{\makebox[0in][l]{\makebox[#1][l]{}\raisebox{\baselineskip}[0in][0in]{\raisebox{#2}[0in][0in]{#3}}}}
     \def\rightbox#1{\makebox[0in][r]{#1}}
     \def\centbox#1{\makebox[0in]{#1}}
     \def\topbox#1{\raisebox{-\baselineskip}[0in][0in]{#1}}
     \def\midbox#1{\raisebox{-0.5\baselineskip}[0in][0in]{#1}}

\vspace{3cm}

\title{
	\logo{
Discrete Maths: Maths Olympiad 
	}
}
\author{ G V V Sharma$^{*}$% <-this % stops a space
	\thanks{*The author is with the Department
		of Electrical Engineering, Indian Institute of Technology, Hyderabad
		502285 India e-mail:  gadepall@iith.ac.in. All content in this manual is released under GNU GPL.  Free and open source.}
	
}	
%\title{
%	\logo{Matrix Analysis through Octave}{\begin{center}\includegraphics[scale=.24]{tlc}\end{center}}{}{HAMDSP}
%}


% paper title
% can use linebreaks \\ within to get better formatting as desired
%\title{Matrix Analysis through Octave}
%
%
% author names and IEEE memberships
% note positions of commas and nonbreaking spaces ( ~ ) LaTeX will not break
% a structure at a ~ so this keeps an author's name from being broken across
% two lines.
% use \thanks{} to gain access to the first footnote area
% a separate \thanks must be used for each paragraph as LaTeX2e's \thanks
% was not built to handle multiple paragraphs
%

%\author{<-this % stops a space
%\thanks{}}
%}
% note the % following the last \IEEEmembership and also \thanks - 
% these prevent an unwanted space from occurring between the last author name
% and the end of the author line. i.e., if you had this:
% 
% \author{....lastname \thanks{...} \thanks{...} }
%                     ^------------^------------^----Do not want these spaces!
%
% a space would be appended to the last name and could cause every name on that
% line to be shifted left slightly. This is one of those "LaTeX things". For
% instance, "\textbf{A} \textbf{B}" will typeset as "A B" not "AB". To get
% "AB" then you have to do: "\textbf{A}\textbf{B}"
% \thanks is no different in this regard, so shield the last } of each \thanks
% that ends a line with a % and do not let a space in before the next \thanks.
% Spaces after \IEEEmembership other than the last one are OK (and needed) as
% you are supposed to have spaces between the names. For what it is worth,
% this is a minor point as most people would not even notice if the said evil
% space somehow managed to creep in.



% The paper headers
%\markboth{Journal of \LaTeX\ Class Files,~Vol.~6, No.~1, January~2007}%
%{Shell \MakeLowercase{\textit{et al.}}: Bare Demo of IEEEtran.cls for Journals}
% The only time the second header will appear is for the odd numbered pages
% after the title page when using the twoside option.
% 
% *** Note that you probably will NOT want to include the author's ***
% *** name in the headers of peer review papers.                   ***
% You can use \ifCLASSOPTIONpeerreview for conditional compilation here if
% you desire.




% If you want to put a publisher's ID mark on the page you can do it like
% this:
%\IEEEpubid{0000--0000/00\$00.00~\copyright~2007 IEEE}
% Remember, if you use this you must call \IEEEpubidadjcol in the second
% column for its text to clear the IEEEpubid mark.



% make the title area
\maketitle

\newpage

%\tableofcontents

\bigskip

\renewcommand{\thefigure}{\theenumi}
\renewcommand{\thetable}{\theenumi}
%\renewcommand{\theequation}{\theenumi}

%\begin{abstract}
%%\boldmath
%In this letter, an algorithm for evaluating the exact analytical bit error rate  (BER)  for the piecewise linear (PL) combiner for  multiple relays is presented. Previous results were available only for upto three relays. The algorithm is unique in the sense that  the actual mathematical expressions, that are prohibitively large, need not be explicitly obtained. The diversity gain due to multiple relays is shown through plots of the analytical BER, well supported by simulations. 
%
%\end{abstract}
% IEEEtran.cls defaults to using nonbold math in the Abstract.
% This preserves the distinction between vectors and scalars. However,
% if the journal you are submitting to favors bold math in the abstract,
% then you can use LaTeX's standard command \boldmath at the very start
% of the abstract to achieve this. Many IEEE journals frown on math
% in the abstract anyway.

% Note that keywords are not normally used for peerreview papers.
%\begin{IEEEkeywords}
%Cooperative diversity, decode and forward, piecewise linear
%\end{IEEEkeywords}



% For peer review papers, you can put extra information on the cover
% page as needed:
% \ifCLASSOPTIONpeerreview
% \begin{center} \bfseries EDICS Category: 3-BBND \end{center}
% \fi
%
% For peerreview papers, this IEEEtran command inserts a page break and
% creates the second title. It will be ignored for other modes.
%\IEEEpeerreviewmaketitle

\begin{abstract}
This book provides a collection of the international maths olympiad problems in discrete mathematics.
\end{abstract}
%Download python codes using 
%\begin{lstlisting}
%svn co https://github.com/gadepall/school/trunk/ncert/computation/codes
%\end{lstlisting}

\renewcommand{\theequation}{\theenumi}
\begin{enumerate}[label=\arabic*.,ref=\theenumi]
%\begin{enumerate}[label=\arabic*.,ref=\thesubsection.\theenumi]
\numberwithin{equation}{enumi}
\documentclass[journal,12pt,twocolumn]{IEEEtran}

\usepackage{amsthm}
\usepackage{mathrsfs}
\usepackage{txfonts}
\usepackage{stfloats}
\usepackage{cite}
\usepackage{cases}
\usepackage{mathtools}
\usepackage{caption}
\usepackage{enumerate}	
\usepackage{enumitem}
\usepackage{amsmath}
%\usepackage{xtab}
\usepackage{longtable}
\usepackage{multirow}
%\usepackage{algorithm}
%\usepackage{algpseudocode}
\usepackage{enumitem}
\usepackage{mathtools}
\usepackage{hyperref}
%\usepackage[framemethod=tikz]{mdframed}
\usepackage{listings}
    %\usepackage[latin1]{inputenc}                                 %%
    \usepackage{color}                                            %%
    \usepackage{array}                                            %%
    \usepackage{longtable}                                        %%
    \usepackage{calc}                                             %%
    \usepackage{multirow}                                         %%
    \usepackage{hhline}                                           %%
    \usepackage{ifthen}                                           %%
  %optionally (for landscape tables embedded in another document): %%
    \usepackage{lscape}     


\usepackage{url}
\def\UrlBreaks{\do\/\do-}


%\usepackage{stmaryrd}


%\usepackage{wasysym}
%\newcounter{MYtempeqncnt}
\DeclareMathOperator*{\Res}{Res}
%\renewcommand{\baselinestretch}{2}
\renewcommand\thesection{\arabic{section}}
\renewcommand\thesubsection{\thesection.\arabic{subsection}}
\renewcommand\thesubsubsection{\thesubsection.\arabic{subsubsection}}

\renewcommand\thesectiondis{\arabic{section}}
\renewcommand\thesubsectiondis{\thesectiondis.\arabic{subsection}}
\renewcommand\thesubsubsectiondis{\thesubsectiondis.\arabic{subsubsection}}

% correct bad hyphenation here
\hyphenation{op-tical net-works semi-conduc-tor}

%\lstset{
%language=C,
%frame=single, 
%breaklines=true
%}

%\lstset{
	%%basicstyle=\small\ttfamily\bfseries,
	%%numberstyle=\small\ttfamily,
	%language=Octave,
	%backgroundcolor=\color{white},
	%%frame=single,
	%%keywordstyle=\bfseries,
	%%breaklines=true,
	%%showstringspaces=false,
	%%xleftmargin=-10mm,
	%%aboveskip=-1mm,
	%%belowskip=0mm
%}

%\surroundwithmdframed[width=\columnwidth]{lstlisting}
\def\inputGnumericTable{}                                 %%
\lstset{
%language=C,
frame=single, 
breaklines=true,
columns=fullflexible
}
 

\begin{document}
%


\theoremstyle{definition}
\newtheorem{theorem}{Theorem}[section]
\newtheorem{problem}{Problem}
\newtheorem{proposition}{Proposition}[section]
\newtheorem{lemma}{Lemma}[section]
\newtheorem{corollary}[theorem]{Corollary}
\newtheorem{example}{Example}[section]
\newtheorem{definition}{Definition}[section]
%\newtheorem{algorithm}{Algorithm}[section]
%\newtheorem{cor}{Corollary}
\newcommand{\BEQA}{\begin{eqnarray}}
\newcommand{\EEQA}{\end{eqnarray}}
\newcommand{\define}{\stackrel{\triangle}{=}}

\bibliographystyle{IEEEtran}
%\bibliographystyle{ieeetr}

\providecommand{\nCr}[2]{\,^{#1}C_{#2}} % nCr
\providecommand{\nPr}[2]{\,^{#1}P_{#2}} % nPr
\providecommand{\mbf}{\mathbf}
\providecommand{\pr}[1]{\ensuremath{\Pr\left(#1\right)}}
\providecommand{\qfunc}[1]{\ensuremath{Q\left(#1\right)}}
\providecommand{\sbrak}[1]{\ensuremath{{}\left[#1\right]}}
\providecommand{\lsbrak}[1]{\ensuremath{{}\left[#1\right.}}
\providecommand{\rsbrak}[1]{\ensuremath{{}\left.#1\right]}}
\providecommand{\brak}[1]{\ensuremath{\left(#1\right)}}
\providecommand{\lbrak}[1]{\ensuremath{\left(#1\right.}}
\providecommand{\rbrak}[1]{\ensuremath{\left.#1\right)}}
\providecommand{\cbrak}[1]{\ensuremath{\left\{#1\right\}}}
\providecommand{\lcbrak}[1]{\ensuremath{\left\{#1\right.}}
\providecommand{\rcbrak}[1]{\ensuremath{\left.#1\right\}}}
\theoremstyle{remark}
\newtheorem{rem}{Remark}
\newcommand{\sgn}{\mathop{\mathrm{sgn}}}
\providecommand{\abs}[1]{\left\vert#1\right\vert}
\providecommand{\res}[1]{\Res\displaylimits_{#1}} 
\providecommand{\norm}[1]{\lVert#1\rVert}
\providecommand{\mtx}[1]{\mathbf{#1}}
\providecommand{\mean}[1]{E\left[ #1 \right]}
\providecommand{\fourier}{\overset{\mathcal{F}}{ \rightleftharpoons}}
%\providecommand{\hilbert}{\overset{\mathcal{H}}{ \rightleftharpoons}}
\providecommand{\system}{\overset{\mathcal{H}}{ \longleftrightarrow}}
	%\newcommand{\solution}[2]{\textbf{Solution:}{#1}}
\newcommand{\solution}{\noindent \textbf{Solution: }}
\newcommand{\myvec}[1]{\ensuremath{\begin{pmatrix}#1\end{pmatrix}}}
\providecommand{\dec}[2]{\ensuremath{\overset{#1}{\underset{#2}{\gtrless}}}}
\DeclarePairedDelimiter{\ceil}{\lceil}{\rceil}
%\numberwithin{equation}{section}
%\numberwithin{problem}{subsection}
%\numberwithin{definition}{subsection}
\makeatletter
\@addtoreset{figure}{section}
\makeatother

\let\StandardTheFigure\thefigure
%\renewcommand{\thefigure}{\theproblem.\arabic{figure}}
\renewcommand{\thefigure}{\thesection}


%\numberwithin{figure}{subsection}

%\numberwithin{equation}{subsection}
%\numberwithin{equation}{section}
%\numberwithin{equation}{problem}
%\numberwithin{problem}{subsection}
\numberwithin{problem}{section}
%%\numberwithin{definition}{subsection}
%\makeatletter
%\@addtoreset{figure}{problem}
%\makeatother
\makeatletter
\@addtoreset{table}{section}
\makeatother

\let\StandardTheFigure\thefigure
\let\StandardTheTable\thetable
\let\vec\mathbf
%%\renewcommand{\thefigure}{\theproblem.\arabic{figure}}
%\renewcommand{\thefigure}{\theproblem}

%%\numberwithin{figure}{section}

%%\numberwithin{figure}{subsection}



\def\putbox#1#2#3{\makebox[0in][l]{\makebox[#1][l]{}\raisebox{\baselineskip}[0in][0in]{\raisebox{#2}[0in][0in]{#3}}}}
     \def\rightbox#1{\makebox[0in][r]{#1}}
     \def\centbox#1{\makebox[0in]{#1}}
     \def\topbox#1{\raisebox{-\baselineskip}[0in][0in]{#1}}
     \def\midbox#1{\raisebox{-0.5\baselineskip}[0in][0in]{#1}}

\vspace{3cm}

\title{ 
	\logo{Discrete}
}
\maketitle
\renewcommand{\thefigure}{\theenumi}
\renewcommand{\thetable}{\theenumi}



\textbf {\large}

\begin{enumerate}[label=\arabic*]
\numberwithin{equation}{enumi}

\item Prove that the fraction $\frac{21n+4}{14n+3}$ is irreducible for every natural number n.







\end{document} 
\item Consider triangle $P_1P_2P_3$ and a point P within the triangle. Lines $P_1P$, $P_2P$, $P_3P$ intersect the opposite sides in points $Q_1$, $Q_2$, $Q_3$ respectively. Prove that, of the numbers\\
\\ $\frac{P_1P}{PQ_1}$,$\frac{P_2P}{PQ_2}$,$\frac{P_3P}{PQ_3}$\\
\\ at least one is $\leq 2$ and at least one is $\geq 2$.\\

\item Construct triangle ABC if AC = b, AB = c and $\angle AMB = \omega $, where M is the midpoint of segment BC and $\omega < 90^o$. Prove that a solution exists if and only if\\
\\ $btan\frac{\omega}{2}\leq c < b$.\\
\\ In what case does the equality hold?\\

\item Consider a plane $\varepsilon$ and three non-collinear points A, B, C on the same side of $\varepsilon$; suppose the plane determined by these three points is not parallel to $\varepsilon$. In plane a take three arbitrary points $A_0, B_0, C_0$. Let L, M, N be the midpoints of segments $AA^{'}$, $BB^{'}$, $CC^{'}$; let G be the centroid of triangle LMN. (We will not consider positions of the points $A^{'}$, $B^{'}$, $C^{'}$ such that the points L, M, N do not form a triangle.) What is the locus of point G as $A^{'}$, $B^{'}$, $C^{'}$ range independently over the plane $\varepsilon$?

 

\item Determine all real numbers x which satisfy the inequality:
\begin{align*}
\sqrt{3-x}-\sqrt{x+1}>\frac{1}{2}
\end{align*}

\item Solve the equation
\begin{align*}
\cos^{2}x + \cos^{2}2x + \cos^{2}3x = 1
\end{align*}

\documentclass[journal,12pt,twocolumn]{IEEEtran}
\usepackage{setspace}
\usepackage{gensymb}
\usepackage{caption}
%\usepackage{multirow}
%\usepackage{multicolumn}
%\usepackage{subcaption}
%\doublespacing
\singlespacing
\usepackage{csvsimple}
\usepackage{amsmath}
\usepackage{multicol}
%\usepackage{enumerate}
\usepackage{amssymb}
%\usepackage{graphicx}
\usepackage{newfloat}
%\usepackage{syntax}
\usepackage{listings}
%\usepackage{iithtlc}
\usepackage{color}
\usepackage{tikz}
\usetikzlibrary{shapes,arrows}



%\usepackage{graphicx}
%\usepackage{amssymb}
%\usepackage{relsize}
%\usepackage[cmex10]{amsmath}
%\usepackage{mathtools}
%\usepackage{amsthm}
%\interdisplaylinepenalty=2500
%\savesymbol{iint}
%\usepackage{txfonts}
%\restoresymbol{TXF}{iint}
%\usepackage{wasysym}
\usepackage{amsthm}
\usepackage{mathrsfs}
\usepackage{txfonts}
\usepackage{stfloats}
\usepackage{cite}
\usepackage{cases}
\usepackage{mathtools}
\usepackage{caption}
\usepackage{enumerate}
\usepackage{tfrupee}	
\usepackage{enumitem}
\usepackage{amsmath}
%\usepackage{xtab}
\usepackage{longtable}
\usepackage{multirow}
%\usepackage{algorithm}
%\usepackage{algpseudocode}
\usepackage{enumitem}
\usepackage{gensymb}
\usepackage{mathtools}
\usepackage{hyperref}
%\usepackage[framemethod=tikz]{mdframed}
\usepackage{listings}
    %\usepackage[latin1]{inputenc}                                 %%
    \usepackage{color}                                            %%
    \usepackage{array}                                            %%
    \usepackage{longtable}                                        %%
    \usepackage{calc}                                             %%
    \usepackage{multirow}                                         %%
    \usepackage{hhline}                                           %%
    \usepackage{ifthen}                                           %%
  %optionally (for landscape tables embedded in another document): %%
    \usepackage{lscape}     


\usepackage{url}
\def\UrlBreaks{\do\/\do-}


%\usepackage{stmaryrd}


%\usepackage{wasysym}
%\newcounter{MYtempeqncnt}
\DeclareMathOperator*{\Res}{Res}
%\renewcommand{\baselinestretch}{2}
\renewcommand\thesection{\arabic{section}}
\renewcommand\thesubsection{\thesection.\arabic{subsection}}
\renewcommand\thesubsubsection{\thesubsection.\arabic{subsubsection}}

\renewcommand\thesectiondis{\arabic{section}}
\renewcommand\thesubsectiondis{\thesectiondis.\arabic{subsection}}
\renewcommand\thesubsubsectiondis{\thesubsectiondis.\arabic{subsubsection}}

% correct bad hyphenation here
\hyphenation{op-tical net-works semi-conduc-tor}

%\lstset{
%language=C,
%frame=single, 
%breaklines=true
%}

%\lstset{
	%%basicstyle=\small\ttfamily\bfseries,
	%%numberstyle=\small\ttfamily,
	%language=Octave,
	%backgroundcolor=\color{white},
	%%frame=single,
	%%keywordstyle=\bfseries,
	%%breaklines=true,
	%%showstringspaces=false,
	%%xleftmargin=-10mm,
	%%aboveskip=-1mm,
	%%belowskip=0mm
%}

%\surroundwithmdframed[width=\columnwidth]{lstlisting}
\def\inputGnumericTable{}                                 %%
\lstset{
%language=C,
frame=single, 
breaklines=true,
columns=fullflexible
}
 

\begin{document}
%
\tikzstyle{block} = [rectangle, draw,
    text width=3em, text centered, minimum height=3em]
\tikzstyle{sum} = [draw, circle, node distance=3cm]
\tikzstyle{input} = [coordinate]
\tikzstyle{output} = [coordinate]
\tikzstyle{pinstyle} = [pin edge={to-,thin,black}]

\theoremstyle{definition}
\newtheorem{theorem}{Theorem}[section]
\newtheorem{problem}{Problem}
\newtheorem{proposition}{Proposition}[section]
\newtheorem{lemma}{Lemma}[section]
\newtheorem{corollary}[theorem]{Corollary}
\newtheorem{example}{Example}[section]
\newtheorem{definition}{Definition}[section]
%\newtheorem{algorithm}{Algorithm}[section]
%\newtheorem{cor}{Corollary}
\newcommand{\BEQA}{\begin{eqnarray}}
\newcommand{\EEQA}{\end{eqnarray}}
\newcommand{\define}{\stackrel{\triangle}{=}}

\bibliographystyle{IEEEtran}
%\bibliographystyle{ieeetr}

\providecommand{\nCr}[2]{\,^{#1}C_{#2}} % nCr
\providecommand{\nPr}[2]{\,^{#1}P_{#2}} % nPr
\providecommand{\mbf}{\mathbf}
\providecommand{\pr}[1]{\ensuremath{\Pr\left(#1\right)}}
\providecommand{\qfunc}[1]{\ensuremath{Q\left(#1\right)}}
\providecommand{\sbrak}[1]{\ensuremath{{}\left[#1\right]}}
\providecommand{\lsbrak}[1]{\ensuremath{{}\left[#1\right.}}
\providecommand{\rsbrak}[1]{\ensuremath{{}\left.#1\right]}}
\providecommand{\brak}[1]{\ensuremath{\left(#1\right)}}
\providecommand{\lbrak}[1]{\ensuremath{\left(#1\right.}}
\providecommand{\rbrak}[1]{\ensuremath{\left.#1\right)}}
\providecommand{\cbrak}[1]{\ensuremath{\left\{#1\right\}}}
\providecommand{\lcbrak}[1]{\ensuremath{\left\{#1\right.}}
\providecommand{\rcbrak}[1]{\ensuremath{\left.#1\right\}}}
\theoremstyle{remark}
\newtheorem{rem}{Remark}
\newcommand{\sgn}{\mathop{\mathrm{sgn}}}
\providecommand{\abs}[1]{\left\vert#1\right\vert}
\providecommand{\res}[1]{\Res\displaylimits_{#1}} 
\providecommand{\norm}[1]{\left\Vert#1\right\Vert}
\providecommand{\mtx}[1]{\mathbf{#1}}
\providecommand{\mean}[1]{E\left[ #1 \right]}
\providecommand{\fourier}{\overset{\mathcal{F}}{ \rightleftharpoons}}
%\providecommand{\hilbert}{\overset{\mathcal{H}}{ \rightleftharpoons}}
\providecommand{\system}{\overset{\mathcal{H}}{ \longleftrightarrow}}
	%\newcommand{\solution}[2]{\textbf{Solution:}{#1}}
\newcommand{\solution}{\noindent \textbf{Solution: }}
\newcommand{\myvec}[1]{\ensuremath{\begin{pmatrix}#1\end{pmatrix}}}
\providecommand{\dec}[2]{\ensuremath{\overset{#1}{\underset{#2}{\gtrless}}}}
\DeclarePairedDelimiter{\ceil}{\lceil}{\rceil}
%\numberwithin{equation}{section}
%\numberwithin{problem}{subsection}
%\numberwithin{definition}{subsection}
\makeatletter
\@addtoreset{figure}{section}
\makeatother

\let\StandardTheFigure\thefigure
%\renewcommand{\thefigure}{\theproblem.\arabic{figure}}
\renewcommand{\thefigure}{\thesection}


%\numberwithin{figure}{subsection}

%\numberwithin{equation}{subsection}
%\numberwithin{equation}{section}
%\numberwithin{equation}{problem}
%\numberwithin{problem}{subsection}
\numberwithin{problem}{section}
%%\numberwithin{definition}{subsection}
%\makeatletter
%\@addtoreset{figure}{problem}
%\makeatother
\makeatletter
\@addtoreset{table}{section}
\makeatother

\let\StandardTheFigure\thefigure
\let\StandardTheTable\thetable
\let\vec\mathbf
%%\renewcommand{\thefigure}{\theproblem.\arabic{figure}}
%\renewcommand{\thefigure}{\theproblem}

%%\numberwithin{figure}{section}

%%\numberwithin{figure}{subsection}



\def\putbox#1#2#3{\makebox[0in][l]{\makebox[#1][l]{}\raisebox{\baselineskip}[0in][0in]{\raisebox{#2}[0in][0in]{#3}}}}
     \def\rightbox#1{\makebox[0in][r]{#1}}
     \def\centbox#1{\makebox[0in]{#1}}
     \def\topbox#1{\raisebox{-\baselineskip}[0in][0in]{#1}}
     \def\midbox#1{\raisebox{-0.5\baselineskip}[0in][0in]{#1}}

\vspace{3cm}

\title{ 
%	\logo{
Discrete
%	}
}

\author{ G V V Sharma$^{*}$% <-this % stops a space
	\thanks{*The author is with the Department
		of Electrical Engineering, Indian Institute of Technology, Hyderabad
		502285 India e-mail:  gadepall@iith.ac.in. All content in this manual is released under GNU GPL.  Free and open source.}
	
}	

\maketitle

%\tableofcontents

\bigskip

\renewcommand{\thefigure}{\theenumi}
\renewcommand{\thetable}{\theenumi}



\begin{enumerate}[label=\arabic*]
\numberwithin{equation}{enumi}
\item In an n-gon all of whose interior angles are equal, the lengths of consecutive
sides satisfy the relation\\
\\$a_1\geq a_2 \geq...... \geq a_n$\\
\\Prove that $a_1 = a_2 = ......= a_n.$\\
\end{enumerate}
\end{document}
     
\item 
\begin{enumerate}
\item Find all positive integers n for which $2^n$ - 1 is divisible by 7.
\item Prove that there is no positive integer n for which $2^n$ + 1 is divisible by 7.
\end{enumerate}

 
\item Given the tetrahedron ABCD whose edges AB and CD have lengths a and b respectively. The distance between the skew lines AB and CD is d, and the angle between them is $\omega$. Tetrahedron ABCD is divided into two solids by plane $\varepsilon$, parallel to lines AB and CD. The ratio of the distances of $\varepsilon$ from AB and CD is equal to k. Compute the ratio of the volumes of the two solids obtained.
\item Consider a triangle OAB with acute angle AOB. Through a point $M \neq O$ perpendiculars are drawn to OA and OB, the feet of which are P and Q respectively. The point of intersection of the altitudes of $\triangle{OPQ}$ is H. What is the locus of H if M is permitted to range over
\begin{enumerate}
\item the side AB, 
\item  the interior of $\triangle{OAB}$?
\end{enumerate}
\item In a plane a set of n points ($n \geq 3$) is given. Each pair of points is connected by a segment. Let d be the length of the longest of these segments. We define a diameter of the set to be any connecting segment of length d. Prove that the number of diameters of the given set is at most n.



 
\item Prove that for every natural number n, and for every real number x$\neq \frac{k\pi}{2^t}$(t = 0,1,..., n; k any integer)
\begin{align*}
\frac{1}{\sin{2x}} + \frac{1}{\sin{4x}}+....\frac{1}{\sin{2^{n}x}} = \cot{x} - \cot{2^{n}x}.
\end{align*}


 
\item In a sports contest, there were m medals awarded on n successive days $(n>1)$. On the first day, one medal and $\frac{1}{7}$ of the remaining m-1 medals were awarded. On the second day, two medals and $\frac{1}{7}$ of the now remaining medals were awarded; and so on. On the $n^{th}$ and last day, the remaining n medals were awarded. How many days did the contest last, and how many medals were awarded altogether?

 
\item Let f be a real-valued function defined for all real numbers x such that, for some positive constant a, the equation
\begin{align}
f(x + a) = \frac{1}{2} + \sqrt{f(x) - f(x)^2}
\end{align}
holds for all x.
\begin{enumerate}
\item Prove that the function f is periodic (i.e., there exists a positive number b such that f(x + b) = f(x) for all x).
\item For a = 1, give an example of a non-constant function with the required properties.
\end{enumerate}

\item For every natural number n, evaluate the sum\\
\\$\sum_{k=1}^{\infty}[\frac{n+2^{k}}{2^{k+1}}] = [\frac{n+1}{2}] + [\frac{n+2}{4}]...........[\frac{n+2^{k}}{2^{k+1}}]$+........\\
\\(The symbol [x] denotes the greatest integer not exceeding x.)


 
\item For each value of k = 1, 2, 3, 4, 5, find necessary and sufficient conditions on the number $a > 0$ so that there exists a tetrahedron with k edges of length a, and the remaining 6-k edges of length 1.\\
\item Let $a_1, a_2,......, a_n$ be real constants, x a real variable, and\\
\\$f(x) = \cos(a_1 + x) + \frac{1}{2}\cos(a_2 + x) + \frac{1}{4}\cos(a_3 + x) +.......+\frac{1}{2^{n} - 1}\cos(a_n + x).$\\
\\Given that $f(x_1) = f(x_2) = 0,$ prove that $x_2 -  x_1 = m\pi$ for some integer m.

 



\item Let a, b and n be integers greater than 1, and let a and b be the bases of two number systems. $A_{n-1}$ and $A_n$ are numbers in the system with base a,  and $B_{n-1}$ and $B_n$ are numbers in the system with base b; these are related as follows:\\
\\ $A_n = x_nx_{n-1}....... x_0, A_{n-1} = x_{n-1}x_{n-2}........x_0,$\\
\\ $B_n = x_nx_{n-1}....... x_0, B_{n-1} = x_{n-1}x_{n-2}........x_0,$\\
\\ $x_n \neq 0,  x_{n-1} \neq 0$\\
\\ Prove:\\
$\frac{A_{n-1}}{A_n} < \frac{B_{n-1}}{B_n}$ if and only if $a > b$\\

\item The real numbers $a_0, a_1,......, a_n,$...... satisfy the condition:\\
\\ 1 = $a_0 \leq a_1 \leq a_2 \leq ..... \leq a_n \leq ........$\\
\\The numbers $b_1, b_2,......, b_n,$....... are defined by\\
\\ $b_n = \sum_{k=1}^{n}(1-\frac{a_{k-1}}{a_k})(\frac{1}{\sqrt{a_k}})$\\
\begin{enumerate}
\item Prove that 0 $\leq bn < 2$ for all n.
\item Given c with 0 $\leq c < 2,$ prove that there exist numbers $a_0, a_1,$..... with
the above properties such that $b_n > c$ for large enough n.
\end{enumerate}

 
\item Consider a convex polyhedron $P_1$ with nine vertices $A_1A_2,......, A_9;$ let $P_i$ be the polyhedron obtained from $P_1$ by a translation that moves vertex $A_1$ to $A_i$(i = 2, 3,......, 9). Prove that at least two of the polyhedra $P_1, P_2,...., P_9$ have an interior point in common.
\item All the faces of tetrahedron ABCD are acute-angled triangles. We consider all closed polygonal paths of the form XY ZTX defined as follows: X is a point on edge AB distinct from A and B; similarly, Y, Z, T are interior points of edges BCCD, DA, respectively. Prove:
\begin{enumerate}
\item If $\angle{DAB}+ \angle{BCD} \neq \angle{CDA} + \angle{ABC},$ then among the polygonal paths, there is none of minimal length.
\item If $\angle{DAB}+ \angle{BCD}= \angle{CDA}+ \angle{ABC},$ then there are infinitely many
shortest polygonal paths, their common length being 2AC $\sin{\frac{\alpha}{2}}$ where
$\alpha = \angle{BAC}+\angle{CAD}+\angle{DAB}.$
\end{enumerate}



 
\item Prove that if $n \geq 4,$ every quadrilateral that can be inscribed in a circle can be dissected into n quadrilaterals each of which is inscribable in a circle.
\item Given four distinct parallel planes, prove that there exists a regular tetrahedron with a vertex on each plane.


 
\item A soldier needs to check on the presence of mines in a region having the shape of an equilateral triangle. The radius of action of his detector is equal to half the altitude of the triangle. The soldier leaves from one vertex of the triangle. What path should he follow in order to travel the least possible distance and still accomplish his mission?\\

\item G is a set of non-constant functions of the real variable x of the form f(x) = ax + b, a and b are real numbers, and G has the following properties:
\begin{enumerate}
\item If f and g are in G, then $g \circ f$ is in G; here $(g \circ f)$(x) = g[f(x)].
\item If f is in G, then its inverse $f^{-1}$ is in G; here the inverse of f(x) = ax +b
is $f^{-1}(x) = \frac{(x-b)}{a}.$
\item For every f in G, there exists a real number $x_f$ such that $f(x_f) = x_f.$
Prove that there exists a real number k such that f(k) = k for all f in G.
\end{enumerate}

 
\item Three players A, B and C play the following game: On each of three cards an integer is written. These three numbers p, q, r satisfy $0 < p < q < r.$ The three cards are shuffled and one is dealt to each player. Each then receives the number of counters indicated by the card he holds. Then the cards are shuffled again; the counters remain with the players. This process (shuffling, dealing, giving out counters) takes place for at least two rounds. After the last round, A has 20 counters in all, B has 10 and C has 9. At the last round B received r counters. Who received q counters on the first round?\\

\item Prove that the number $\sum_{k=0}^{n} \begin{pmatrix} 2n+1 \\ 2k+1 \end{pmatrix} 2^{3k}$ is not divisible by 5 for any integer $n \geq 0.$
\item Consider decompositions of an 8 $\times$ 8 chessboard into p non-overlapping rectangles subject to the following conditions:
\item Each rectangle has as many white squares as black squares.
\item  If $a_i$ is the number of white squares in the i-th rectangle, then $a_1 < a_2 <
..... < a_p.$ Find the maximum value of p for which such a decomposition is possible. For this value of p, determine all possible sequences $a_1, a_2,......, a_p.$

 
\item Let $x_i, y_i$ (i = 1, 2,....., n) be real numbers such that
\begin{align*}
x_1 \geq x_2 \geq...... \geq x_n  \text{ and } y_1 \geq y_2 \geq...... \geq y_n.
\end{align*}
Prove that, if $z_1, z_2,...... , z_n$ is any permutation of $y_1, y_2,..... , y_n,$ then\\
$\sum_{i=1}^{n}(x_i - y_i)^2 \leq \sum_{i=1}^{n}(x_i - z_i)^2$
\item Let $a_1, a_2, a_3,.......$ be an infinite increasing sequence of positive integers. Prove that for every p $\geq$ 1 there are infinitely many $a_m$ which can be written in the form $a_m = xa_p + ya_q$ with x, y positive integers and $q > p.$
\item When $4444^{4444}$ is written in decimal notation, the sum of its digits is A. Let
B be the sum of the digits of A. Find the sum of the digits of B. (A and B are written in decimal notation.)


 
\item In a plane convex quadrilateral of area 32, the sum of the lengths of two opposite sides and one diagonal is 16. Determine all possible lengths of the other diagonal.
\item A rectangular box can be filled completely with unit cubes. If one places as many cubes as possible, each with volume 2, in the box, so that their edges are parallel to the edges of the box, one can fill exactly 40$\%$ of the box. Determine the possible dimensions of all such boxes.


 


\item Let a and b be positive integers. When $a^2+b^2$ is divided by a+b, the quotient
is q and the remainder is r. Find all pairs (a, b) such that $q^2 + r = 1977.$\\

\item Four real constants a, b, A, B are given, and\\
\\$f(\theta) = 1 - a \cos\theta  - b \sin\theta - A \cos 2\theta - B \sin 2\theta.$\\
\\Prove that if $f(\theta) \geq 0$ for all real $\theta,$ then $a^2 + b^2 \leq 2$ and $A^2 + B^2 \leq 1.$


 
\item Let $n > 6$ be an integer and $a_1, a_2,....... a_k$ be all the natural numbers less than n and relatively prime to n. If\\
$a_2 - a_1 = a_3 - a_2 =..........= a_k - a_{k-1} > 0$
\\prove that n must be either a prime number or a power of 2.\\

\item An infinite sequence $x_0, x_1, x_2,$....... of real numbers is said to be bounded
if there is a constant C such that $\begin{vmatrix} x_i \end{vmatrix} \leq C$ for every i $\geq 0.$ Given any real number $a > 1,$ construct a bounded infinite sequence $x_0, x_1, x_2,$.......such that\\
\\ $\begin{vmatrix} x_i - x_j \end{vmatrix} \begin{vmatrix} i - j \end{vmatrix}^a \geq 1$\\
\\for every pair of distinct non negative integers i, j.\\


 
\item Let R denote the set of all real numbers. Find all functions f : $R \rightarrow R$
such that\\
\\ $f(x^2+f(y))= y +(f(x))^2$ for all x,y $\in $R.

\item Let S be a finite set of points in three-dimensional space. Let $S_x$, $S_y$, $S_z$ be the sets consisting of the orthogonal projections of the points of S onto the yz-plane, zx-plane, xy-plane, respectively. Prove that\\
\\  $\begin{vmatrix} S^2 \end{vmatrix}$ $\leq$ $\begin{vmatrix} S_x \end{vmatrix}$ . $\begin{vmatrix} S_y \end{vmatrix}$ . $\begin{vmatrix} S_z \end{vmatrix}$\\
\\ where $\begin{vmatrix} A \end{vmatrix}$ denotes the number of elements in the finite set $\begin{vmatrix} A \end{vmatrix}$. (Note: The orthogonal projection of a point onto a plane is the foot of the perpendicular from that point to the plane.)




\item Let D be a point inside acute triangle ABC such that $\angle ADB = \angle ACB  + \frac{\pi}{2}$ and AC · BD = AD · BC.
\begin{enumerate}
\item Calculate the ratio (AB · CD)/(AC · BD).
\item Prove that the tangents at C to the circumcircles of $\triangle ACD$ and $ \triangle BCD$ are perpendicular.
\end{enumerate}



  
\item ABC is an isosceles triangle with AB = AC. Suppose that
\begin{enumerate}
\item M is the midpoint of BC and O is the point on the line AM such that OB is perpendicular to AB;
\item Q is an arbitrary point on the segment BC different from B and C;
\item E lies on the line AB and F lies on the line AC such that E, Q, F are distinct and collinear.
\end{enumerate}



 
\item Let A, B, C, D be four distinct points on a line, in that order. The circles with diameters AC and BD intersect at X and Y . The line XY meets BC at Z. Let P be a point on the line XY other than Z. The line CP intersects the circle with diameter AC at C and M, and the line BP intersects the circle with diameter BD at B and N. Prove that the lines AM, DN, XY are concurrent.\\

\item Let ABCDEF be a convex hexagon with AB = BC = CD and DE = EF = F A, such that $\angle BCD = \angle EFA = \frac{\pi}{3}$. Suppose G and H are points in the interior of the hexagon such that $\angle AGB = \angle DHE = \frac{2\pi}{3}$. Prove that AG + GB + GH + DH + $HE \geq CF$.


 
\item The positive integers a and b are such that the numbers 15a + 16b and 16a - 15b are both squares of positive integers. What is the least possible value that can be taken on by the smaller of these two squares?





 
\item Let $x_1, x_2, .... , x_n$ be real numbers satisfying the conditions\\
\\ $\begin{vmatrix} x_1 + x_2 + .... + x_n \end{vmatrix} = 1$\\
\\ and\\
$\begin{vmatrix}x_i \end{vmatrix} \leq \frac{n+1}{2}$  i = 1, 2, . . . , n\\
\\ Show that there exists a permutation $y_1, y_2, . . . , y_n$ of $x_1, x_2, . . . , x_n$ such that\\
\\ $\begin{vmatrix}y_1+2y_2+...+ny_n \end{vmatrix} \leq \frac{n+1}{2}$\\

\item Find all pairs (a, b) of integers a, $b \geq 1$ that satisfy the equation\\
\\ $a^{b^{2}} = b^a$.\\

\item For each positive integer n , let f(n) denote the number of ways of representing n as a sum of powers of 2 with non-negative integer exponents. Representations which differ only in the ordering of their summands are considered to be the same. For instance, f(4) = 4, because the number 4 can be represented in the
following four ways:\\
\\ 4; 2 + 2; 2 + 1 + 1; 1 + 1 + 1 + 1.\\
\\ Prove that, for any integer $n \geq 3$,\\
\\ $2^\frac{n^2}{4} < f(2^n) < 2^\frac{n^2}{2}$







 

\item For any positive integer n, let d(n) denote the number of positive divisors of n (including 1 and n itself). Determine all positive integers k such that $d(n^2)/d(n)$ = k for some n.\\

\item Determine all pairs (a, b) of positive integers such that $ab^2$ + b + 7 divides $a^2b + a + b.$



 
\item Two circles $G_1$ and $G_2$ are contained inside the circle G, and are tangent to G at the distinct points M and N, respectively. $G_1$ passes through the center of $G_2$. The line passing through the two points of intersection of $G_1$ and $G_2$ meets G at A and B. The lines MA and MB meet $G_1$ at C and D, respectively.

 

\item 100 cards are numbered 1 to 100 (each card different) and placed in 3 boxes (at least one card in each box). How many ways can this be done so that if two boxes are selected and a card is taken from each, then the knowledge of their sum alone is always sufficient to identify the third box?\\

\item Can we find N divisible by just 2000 different primes, so that N divides $2^{N} + 1$? [N may be divisible by a prime power.]


 


\item Prove that\\
\\ $\frac{a}{\sqrt{(a^2 + 8bc)}}$+$\frac{b}{\sqrt{(b^2 + 8ac)}}$+$\frac{c}{\sqrt{(c^2 + 8ba)}}$ $\geq 1$\\
\\ for all positive real numbers a,b and c.\\

\item Let a, b, c, d be integers with $a > b > c > d > 0$. Suppose that\\
\\ ac + bd = (b + d  a - c) (b + d - a + c).\\
\\ Prove that ab + cd is not prime.



 






\item S is the set of all (h, k) with h, k non-negative integers such that h + $k < n$. Each element of S is colored red or blue, so that if (h, k) is red and $h^{'}\leq h, k^{'} \leq k$, then ($h^{'}, k^{'}$) is also red. A type 1 subset of S has n blue elements with different first member and a type 2 subset of S has n blue elements with different second member. Show that there are the same number of type 1 and type 2 subsets.\\

\item Find all pairs of integers $m > 2$, $n > 2$ such that there are infinitely many positive integers k for which $k^n + k^2 - 1$ divides $k^m + k + 1$.\\

\item The positive divisors of the integer $n > 1$ are $d_1 < d_2 < ... < d_k$, so that $d_1 = 1$, $d_k = n$. Let $d = d_1d_2 + d_2d_3 + ... + d_{k-1}d_k$. Show that $d < n^2$ and find all n for which d divides $n^2$.


 
\item A convex hexagon has the property that for any pair of opposite sides the distance between their midpoints is $\sqrt{3}/2$ times the sum of their lengths Show that all the hexagon’s angles are equal.\\

\item ABCD is cyclic. The feet of the perpendicular from D to the lines AB, BC, CA are P, Q, R respectively. Show that the angle bisectors of ABC and CDA meet on the line AC if RP=RQ.

 
\item Consider in plane circle $\Gamma$ with center O and a line l not intersecting circle $\Gamma$. Prove that there is a unique point Q on the perpendicular drawn from O to the line l, such that for any point P on the line l, PQ represents the length of the tangent from P to the circle $\Gamma$.

\item Let ABCD be a quadrilateral X and Y be the mid points of AC and BD respectively and the lines through X and Y respectively parllel to BD, AC meet in O. Let P, Q, R, S be the mid points of AB, BC, CD, DA respectively. Prove that
\begin{enumerate}
\item quadrilaterals APOS and APXS have the same area.
\item the areeas is the quadrilateral APOS, BQOP, CROQ, DSOR are all equal. 
\end{enumerate}
 
\item Let x, y, z be three positive reals such that xyz $ \geq 1$. Prove that
\begin{align*}
\frac{x^5-x^2}{x^5 + y^2 + z^2}+\frac{y^5-y^2}{x^2 + y^5 + z^2}+\frac{z^5-z^2}{x^2 + y^2 + z^5}\geq 0
\end{align*}

\item Let $a_1, a_2,$....... be a sequence of integers with infinitely many positive and negative terms. Suppose that for every positive integer n the
numbers $a_1, a_2,........, a_n$ leave n different remainders upon division by n.
Prove that every integer occurs exactly once in the sequence $a_1, a_2,$.......

\item Determine all positive integers relatively prime to all the terms of the infinite sequence
\begin{align*}
a_n = 2^n + 3^n + 6^{n}-1, n\geq 1.
\end{align*}

\item Find the least possible value of a + b, where a, b are positive integers such that 11 divides a + 13b and 13 divides a+11b.

\item A 6 $\times$ 6 square is dissected in to 9 rectangles by lines parallel to its sides such that all these rectangles have integer sides, Prove that there are always two congruent rectangles.

\item Let ABCD be a quadrilateral in which AB is parallel to CD and perpendicular to AD ; AB = 3CD;and the area of the quadrilateral is 4. If a circle can be drawn touching all the sides of the quadrilateral, find its radius.

\item Let a, b, c be the three natural numbers such that $a < b < c$ and ged(c-a, c-b) = 1. Suppose there exists an integer d such that a+d, b+d, c+d from the sides of a right angled triangle. Prove that there exist integers l, m such that 
c + d = $l^2 + m^2$.

\item Find all pairs (a, b) of real numbers such that whenever $\alpha$ is a root of
\begin{align} 
x^2 + ax + b = 0,
\end{align}
$\alpha^2$ - 2 is also a root of the equation.


\item 
\begin{enumerate}
\item Prove that
\begin{align*}
\frac{x^{2}}{(x - 1)^{2}}+\frac{y^{2}}{(y - 1)^{2}} + \frac{z^{2}}{(z - 1)^{2}}\geq1
\end{align*}
for all real numbers x, y, z, each different from 1, and satisfying xyz = 1.
\item Prove that equality holds above for infinitely many triples of rational numbers x, y, z, each different from 1, and satisfying xyz = 1.
\end{enumerate} 

\item Prove that there exist infinitely many positive integers n such that $n^{2} + 1$ has a prime divisor which is greater than $2n + \sqrt{2n}$.

\item Let n and k be positive integers with k $\geq$ n and k - n an even number. Let 2n lamps labelled 1, 2,....., 2n be given, each of which can be either on or off. Initially all the lamps are off. We consider sequences of steps: at each step one of the lamps is switched (from on to off or from off to on).
Let N be the number of such sequences consisting of k steps and resulting in the state where lamps 1 through n are all on, and lamps n + 1 through 2n are all off.
Let M be the number of such sequences consisting of k steps, resulting in the state where lamps 1 through n are all on, and lamps n + 1 through 2n are all off, but where none of the lamps n + 1 through 2n is ever switched on.
Determine the ratio $\frac{N}{M}$.

\item All the points in the plane are coloured using three colours, Prove that there exists a triangle with vertices having the same colour such that either it is isosceles or its angles are in geometric progression.
 







\item Let 
\begin{align*}
P_1(x) = ax^2 - bx - c,
\end{align*}
\begin{align*}
P_2(x) = bx^2 - cx - a,
\end{align*}
\begin{align*} 
P_3(x) = cx^2 - ax - b
\end{align*}
be three quadratic polynomials where a, b, c are non-zero real numbers. Suppose there exists a real number 
$\alpha$ such that $P_(\alpha) = P_2(\alpha)= P_3(\alpha)$. Prove that a = b = c

\item Find the number of 4-digit numbers  having non-zero digits and which are divisible by 4 but not 8.
 
\item Let ABCD be a quadrilateral inscribed in a circle. Let E, F, G, H be the midpoints of the arcs AB, BC, CD, DA of the circle. Suppose AC.BD = EG.FH. Prove that AC, BD, EG, FH are concurrent.

\item Let D,E,F be points on the sides BC, CA, AB respectively of a triangle ABC such that BD = CE = AF and
 $\angle$BDF = $\angle$CED = $\angle$AFE. Prove that ABC is equilateral.






















\item Let ABC be a triangle and D be a point on the segment BC such that DC = 2BD. Let E be the mid-point of AC. Let AD and BE intersect in P. Determine the ratios BP/PE and AP/PD.

\item Let ABC be a triangle. Let BE and CF be internal angle bisector of $\angle$B and $\angle$C respectively with E on AC and F on AB. Suppose X is a point on the segment CF such that AX $\perp$ CF;and Y is a point on the segment BE such that AY $\perp$ BE. Prove that XY = (b + c - a)/2 where BC = a, CA = b and AB = c.


\item Let AD and BC be the parallel sides of a trapezium ABCD. Let P and Q be the midpoints of the diagonals AC and BD. If $AD = 16$ and $BC = 20$, What is the length of PQ?

\item In a triangle ABC, let H, I and O be the orthocentre, incentre and circumcentre, respectively. If the points B, H, I, C lie on a circle, What is magnitude of $\angle BOC$ in degrees?

\item Let ABC be an equilateral triangle. Let P and S be points on AB and AC,respectively, and let Q and R be points on BC such that PQRS is a triangle. If $PQ = \sqrt{3}PS$ and the area of PQRS is $28\sqrt{3}$, What is the length of PC?

\item Let $A_1, B_1, C_1, D_1$ be the midpoints of the sides of a convex quadrilateral ABCD and let $A_1, B_2, C_2, D_2$ be the midpoints of the quadrilateral $A_1B_1C_1D_1$. If $A_2B_2C_2D_2$ is a rectangle with sides 4 and 6, then what is the product of the lengths of the diagonals of ABCD?

\item Let S be a circle with centre O. A chord AB, not a diameter, divides S into two regions $R_1$ and $R_2$ such that O belongs to $R_2$. Let $S_1$ be a circle with centre in $R_1$, touching AB at X and S internally. Let $S_2$ be a circle with centre in $R_2$, touching AB at Y, the circle S internally and passing through the centre of S. The point X lies on the diameter passing through the centre of  $S_2$ and $\angle YXO = 30^{o}$. If the radius of $S_2$ is 100 then what is the radius of $S_1$?

\item In a triangle ABC with $\angle BCA = 90^{o}$, the perpendicular bisector of AB intersects segments AB and AC at X and Y, respectively. If the ratio of the area of quadrilateral BXYC to the area of triangle ABC is 13 : 18 and BC = 12 then what is the length of AC?

\item The first term of a sequence is 2014. Each succeeding term is the sum of the cubes of the digits of the previous term. What is the $2014^{th}$ term of the sequence?

\item In a triangle with integer side lengths, one side is three times as long as second side, and the length of the third side is 17. What is the greatest possible perimeter of the triangle?

\item For how many natural numbers n between 1 and 2014(both inclusive) is $\frac{8n}{9999 - n}$ an integer?

\item One morning, each member of Manjul's family drank an 8-ounce mixture of coffee and milk. The amounts of coffee and milk varied from cup to cup, but were never zero. Manjul drank 1/7-th of the total amount of milk and 2/17-th of the total amount of coffee. How many people are there in Manjul's family?

\item For a natural number b, let N(b) denote the number of natural numbers a for which the equation $x^2 + ax + b = 0$ has integer roots. What is the smallest value of b for which N(b)=20?

\item Let S be a set of real numbers with mean M. If the means of the sets $S \cup $ \{15\} and $S \cup $ \{15,1\} are 
M + 2 and M + 1 respectively, then how many elements does S have?

\item Let f be a one-to-one funtion from the set of natural numbers to itself such that f(mn) = f(m)f(n) for all natural numbers m and n. What is the least possible value of f(999)?

\item What is the number of ordered pairs (A, B) where A and B are subsets of \{1,2,....,5\} such that neither 
$A \subseteq B$ nor $B \subseteq A$?

\item A man walks a certain distance and rides back in $3\frac{3}{4}$ hours; he could ride both ways in $2\frac{1}{2}$ hours. How many hours would it take him to walk both ways?

\item How many line segments have both their endpoints located at the vertices of a given cube?

\item Let E(n) denote the sum of the even digits of n. For example, E(1243) = 2 + 4 = 6. What is the value of E(1) + E(2) + E(3) +.... + E(100)?

\item How many two-digit positive integers N have the property that the sum of N and the number obtained by reversing the order of the digits of N is a perfect square?

\item A 2 $\times$ 3 rectangle and a 3 $\times$ 4 rectangle are contained within a square without overlapping at any interior point, and the sides of the square are parallel to the sides of the two given rectangles. What is the smallest possible area of the square?

\item What is the greatest possible perimeter of a right-angled triangle with integer side lengths if one of the sides has length 12?

\item At a party, each man danced with exactly four women and each woman danced with exactly three men. Nine men attended the party. How many women attended the party?

\item A subset B of the set of first 100 positive integers has the property that no two elements of B sum to 125. What is the maximum possible number of elements in B?

\item The digits of a positive integer n are four consecutive integers in decreasing order when read from left to right. What is the sum of the possible remainders when n is divided by 37?

\item Let N denote the set of all natural numbers. Define a function T : N $\to$ N by
T(2k) = k and T(2k + 1) = 2k + 2. We write $T^{2}(n)$ = T(T(n)) and in general
$T^{k}(n) = T^{k-1}(T(n))$ for any $k > 1$.
\begin{enumerate}
\item Show that for each n $\in$ N, there exists k such that $T^{k}$(n) = 1.
\item For k $\in$ N, let $c_k$ denote the number of elements in the set $\{n : T^{k}(n) = 1\}$.
Prove that $c_{k+2} = c_{k+1} + c_{k}$s, for $k \geq 1$.
\end{enumerate}

\item Suppose 2016 points of the circumference of a circle are coloured red and the remaining points are coloured blue. Given any natural number n $\geq$ 3, Prove that there is a regular n-sided polygon all of whose vertices are blue. 
\item Consider $n^{2}$ units squares in the xy-plane centred at point (i, j) with integer co-ordinates, 
1 $\leq$ i $\leq$ n, 1 $\leq$ j $\leq$ n. It is required to colour each unit square in such a way that when ever 
1 $\leq$ i $<$ j $\leq$ n and 1 $\leq$ k $<$ l $\leq$ n,the three squares with centres at (i, k),(j, k),(j, l)have distinct colours. What is the least possible number of colours needed?

\item Let x, y, z be real numbers each greater than 1. Prove that 
\begin{align*}
\frac{x+1}{y+1} + \frac{y+1}{z+1} + \frac{z+1}{x+1} \leq \frac{y-1}{z-1} + \frac{z-1}{x-1}.
\end{align*}
 
\item Let D be an interior point of the side BC of a triangle ABC. Let $I_1$ and $I_2$ be the incentres of triangles ABD and ACD respectively. Let $AI_1$ and $AI_2$ meet BC in E and F respectively. If $\angle BI_1E = 60^{o}$, what is the measure of $\angle CI_2F$ in degrees?

\item Let ABC ba an acute-angled triangle and let H be its orthocentre. Let $G_1$, $G_2$ and $G_3$ be the centroids of the triangles HBC, HCA and HAB respectively. If the area of the triangle $G_1G_2G_3$ is 7 units, what is the area of the triangle ABC?

\item Triangles ABC and DEF are such that $\angle A = \angle D$, AB  = DE = 17, BC = EF = 10  and AC - DF = 12. What is AC + DF?

\item In a triangle ABC, right-angled at A, the altitude through A and the internal bisector of $\angle A$ have lengths 3 and 4 respectively. Find the length of the median through A?

\item In a triangle ABC, the median from B to CA is perpendicular to the median from C to AB. If the median from A to BC is 30, determine $(BC^2 + CA^2 + AB^2)/100$.

\item Let AB be a chord of a circle with centre O. Let C be a point on the circle such that $\angle ABC = 30^{o}$ and O lies inside the triangle ABC. Let D be a point on AB such that $\angle DCO = \angle OCB = 20^{o}$. Find the measure of 
$\angle CDO$ in degrees.

\item Let ABCD be a trapezium in which AB$||$ CD and AD $\perp$ AB. Suppose ABCD has  an incircle which touches AB at Q and CD at P. Given that PC = 36 and QB = 49, find PQ.

\item In a quadrilateral ABCD, it is given that AB = AD = 13, BC = CD = 20, BD = 24. If r is the radius of the circle inscribable, then what is the integer closest to r? 

\end{enumerate}

\end{document}


